\section{Solving the boundary system and controlling its energy}\label{four:sec:existence}

We now construct solutions of the system, exhaust the asymptotically flat
end, and compare the resulting Riemannian energy with the energy of the
prepared data.  The distinction between two kinds of constants is essential.
Constants in the scalar solution map and in local regularity may depend
arbitrarily on a fixed $N$; constants used in the final flux estimate are
the polynomial bounds $\Pi_N$ established earlier.

\subsection{A scalar solution map on the whole trial space}

Fix $0<\alpha<1$, a sufficiently large integer $N$, and an allowed
truncation $\Omega_R$.  In this subsection constants may depend on $R$,
$N$, and a bound for the trial input in
\[
 X_R=\{Z\in C^{1,\alpha}(\overline\Omega_R):Z|_{S_R}=0\}.
\]
The four stages of the homotopy are parametrized consecutively by
$s\in[0,4]$, with $s=0$ the desired system and $s=4$ its zero terminal
scale.  Denote their trace prescriptions by
$F_s^{\mathrm{presc}}(x,w,h,t)$.  All dependence suppressed in this
notation is the smooth dependence specified in the homotopy.

\begin{lemma}[Scalar trial problem]\label{four:ex:scalar-trial}
For every $Z\in X_R$ and $s\in[0,4]$, the prescribed trace equation with
the assigned constant Dirichlet value on every boundary component has a
unique solution $f_s[Z]\in C^{3,\alpha}(\overline\Omega_R)$.  The map
$(s,Z)\mapsto f_s[Z]$ is continuous into $C^{3,\alpha}$ and maps bounded
trial sets to bounded sets, uniformly in $s$.  No lower bound for the
implicitly defined $t$ is imposed on the trial inputs.
\end{lemma}

\begin{proof}
We use the scalar method of continuity in the affine space of
$C^{3,\alpha}$ functions with the assigned boundary values, with equation
space $C^{1,\alpha}$. Thus every solution estimated along this
continuation already has the differentiability required by
Theorem~\ref{four:thm:axial-regularity}; the estimates below are uniform
bounds, not an assumption of uniform third derivatives.

The normalized equation is
\begin{equation}\label{four:ex:scalar-normalized}
 \Achi:\Hess_g f=G_s(x,Z,f,\dd f),\qquad
 G_s=\frac{\sqrt D}{l}
       \bigl(F_s^{\mathrm{presc}}-\tr_{\Achi}K\bigr).
\end{equation}
The implicit relation defining $t$ is smoothly solvable for every
$(Z,\dd f)$.  Neither $t$, the matrix, nor the positive prefactor
$\sqrt D/l$ depends on the value of $f$.  The prescriptions satisfy
\[
 |F_s^{\mathrm{presc}}-h|\le C,
 \qquad \partial_h F_s^{\mathrm{presc}}\ge1,
 \qquad h=\eta_N f,
\]
with the other arguments held fixed.  For a second parameter
$\beta\in[0,1]$, use the scalar interpolation
\begin{equation}\label{four:ex:scalar-interpolation}
 A_\beta:\Hess_g f=G_\beta,
 \quad A_\beta=(1-\beta)\Id+\beta\Achi,
 \quad G_\beta=(1-\beta)\eta_N f+\beta G_s,
\end{equation}
with the same Dirichlet values.  At a positive interior maximum above a
fixed height threshold, both source summands are positive, whereas the
left side is nonpositive.  The negative minimum argument is identical
with signs reversed.  Including the assigned boundary values gives
$|h|\le C$, uniformly in both parameters and the bounded trial set.

We next check the degeneracy at large gradient, before invoking any
regularity theorem.  Put $\sigma=|\dd f|_g$.  Uniformly for bounded $Z$,
$\sigma\to\infty$ forces $t\to-\infty$: a lower bound on $t$ would bound
$e^{6t}l(t)^2$ away from zero, contradicting
$e^{6Z}=e^{6t}(1+l(t)^2\sigma^2)$.  Thus eventually $t<0$ and
$l(t)=e^{Nt}$, giving the exact equation
\begin{equation}\label{four:ex:large-gradient-equation}
 e^{6Z}=e^{6t}+e^{(2N+6)t}\sigma^2.
\end{equation}
Writing $\gamma_N=6/(N+3)$, its consequences, uniformly for $Z$ in a
fixed bounded interval, are
\begin{align}
 t&=-\frac{\log\sigma}{N+3}+\frac{3Z}{N+3}+o(1),\nonumber\\
 d&=e^{-6NZ/(N+3)}\sigma^{-\gamma_N}(1+o(1)),
 &\chi&=\frac{3}{N+3}d(1+o(1)),\label{four:ex:large-gradient}\\
 \frac{\sqrt D}{l}&=\sigma(1+o(1)).\nonumber
\end{align}
For example, the first line follows by taking logarithms of
$e^{(2N+6)t}\sigma^2=e^{6Z}(1-d)$; the last follows directly from
$\sqrt D/l=(\sigma^2+l^{-2})^{1/2}$.
Since $N\ge4$, one has $\gamma_N<1$.  Compactness of the remaining
bounded-gradient range therefore gives constants $c,C>0$ such that
\begin{equation}\label{four:ex:linear-growth}
 \chi_\beta:=1-\beta+\beta\chi\ge\frac{c}{1+\sigma},
 \qquad |G_\beta|\le C(1+\sigma).
\end{equation}
The estimate for the source uses the height bound just proved.  Notice
also the exact structural identity
\[
 A_\beta=\Id-(1-\chi_\beta)e\otimes e,
 \qquad e=\grad f/|\grad f|,
 \qquad 0<\chi_\beta\le1.
\]
In particular, all three transverse eigenvalues remain exactly one
throughout the scalar interpolation.

Here is a gradient estimate using only \eqref{four:ex:linear-growth}.
Extend $g$ smoothly across the boundary to a compact enlargement.  All
distances below are distances of this extension, at scales below its
injectivity radius.  If $H_0$ is the uniform bound for $|f|$, set
\[
 \psi(r)=\frac1k\log(1+B'r),\qquad
 q_1(r)=\psi'(r),\qquad \psi''(r)=-kq_1(r)^2.
\]
First choose $k$ large relative to the constants in
\eqref{four:ex:linear-growth} and the bounded geometry.  Next choose
$h_0<1/(4k)$ below all collar, injectivity, and distinct-boundary-component
separation scales.  Finally choose $B'$ so large that
\begin{equation}\label{four:ex:log-choices}
 q_1\ge1\quad(0\le r\le h_0),\qquad
 \psi(h_0)>2H_0.
\end{equation}
These requirements are compatible: $q_1(h_0)$ tends to $1/(kh_0)>4$
and $\psi(h_0)$ tends to infinity as $B'\to\infty$.

On the inward collar of a boundary face with constant value $f_0$,
let $r$ be its distance function.  The functions
$f_0\pm\psi(r)$ are upper and lower barriers.  To check the upper
barrier at a hypothetical positive maximum of
$f-f_0-\psi(r)$, the common gradient is $q_1\grad r$.  The tangential
Hessian of the barrier contributes at most $Cq_1$, and its axial
contribution is at most $-ckq_1^2/(1+q_1)$.  By enlarging $k$ this is
strictly less than $-C(1+q_1)$, the lower bound for the equation's
source at the solution height.  Ellipticity and the second derivative
test contradict the equation.  For the lower barrier both inequalities
reverse.  At the outer edge of the collar the comparisons follow from
\eqref{four:ex:log-choices}.  Consequently
\begin{equation}\label{four:ex:boundary-modulus}
 |f(x)-f_0|\le\psi\bigl(\operatorname{dist}(x,\mathrm{face})\bigr)
 \quad\hbox{in its $h_0$ collar}.
\end{equation}

For completeness, the same modulus controls interior pairs.  Suppose
that
\[
 f(x)-f(y)-\psi\bigl(\operatorname{dist}(x,y)\bigr)
\]
has a positive maximum on pairs of distance at most $h_0$.  Neither
the diagonal nor the distance-$h_0$ set contains such a maximum.
If an endpoint lies on a face, the other endpoint lies in that face's
collar; \eqref{four:ex:boundary-modulus} and monotonicity of $\psi$ exclude
this as well.  Thus both endpoints are interior.  Let $r>0$ be their
distance and orient their unique short geodesic from $y$ to $x$.
The first derivative test says that their gradients are $q_1(r)$ times
the corresponding parallel unit tangent vectors.

Vary both endpoints in each of three parallel transverse directions.
The index-form estimate for the second variation of distance, using
the parallel field along the geodesic, is bounded above by $Cr$.
The second derivative test for the two-point maximum therefore gives
\[
 \tr_{e^\perp}\Hess f(x)-\tr_{e^\perp}\Hess f(y)
 \le Cq_1r.
\]
Separate axial variations at the two endpoints give
\[
 \Hess f(x)(e,e)\le\psi'',\qquad
 \Hess f(y)(e,e)\ge-\psi''.
\]
Because the transverse coefficients are exactly one at both endpoints,
while the axial coefficients satisfy \eqref{four:ex:linear-growth}, the
difference of the two left sides of \eqref{four:ex:scalar-interpolation} is
at most
\[
 Cq_1r-\frac{2ckq_1^2}{1+q_1}.
\]
Their source difference is at least $-2C(1+q_1)$.  Since $q_1\ge1$,
these inequalities contradict each other for the already sufficiently
large $k$.  No comparison of the two axial coefficients was used.
Only variations of the interior endpoints were required, so the
short geodesic need not remain inside $\Omega_R$.
We have proved the two-point modulus; letting $y\to x$ gives
$|\dd f|\le\psi'(0)=B'/k$ everywhere, including the boundary by
\eqref{four:ex:boundary-modulus}.

The scalar interpolation is now uniformly elliptic.  On the given
smooth interior and constant-Dirichlet boundary patches, its equation
has the exact form required by Theorem~\ref{four:thm:axial-regularity}:
the gradient and source are bounded, and
$0<\chi_0\le\chi_\beta\le1$.  That theorem first gives a uniform
H\"older bound for $\dd f$, at some positive exponent.  The coefficients
and source in \eqref{four:ex:scalar-interpolation} are smooth compositions
of $x,Z,f,\dd f$, so the scalar Dirichlet Schauder estimate gives
$C^{2,\alpha'}$ bounds for some $\alpha'>0$.  These bounds make
$f,\dd f$ Lipschitz.  Together with $Z\in C^{1,\alpha}$, this yields
$C^{0,\alpha}$ coefficients and source, hence $C^{2,\alpha}$ bounds.
The compositions are now $C^{1,\alpha}$ and the next Dirichlet estimate
gives a $C^{3,\alpha}$ bound.

Subtract a smooth extension of the assigned boundary values and work
in the affine space with model
$Y_R=C^{3,\alpha}_0(\overline\Omega_R)$, where the zero trace is on the
entire boundary.  The equation defines a continuously differentiable
map into $C^{1,\alpha}(\overline\Omega_R)$.  Its derivative in $f$ is
\[
 L\varphi=A_\beta^{ij}\nabla_i\nabla_j\varphi
          +b^i\nabla_i\varphi-c(x)\varphi,
 \qquad
 c(x)=\eta_N\left[(1-\beta)+
 \beta\frac{\sqrt D}{l}\partial_hF_s^{\mathrm{presc}}\right]>0.
\]
Derivatives of $A_\beta$ and $t$ with respect to $\dd f$ contribute
only first derivative terms in $\varphi$.  Its coefficients are
$C^{1,\alpha}$ and the ellipticity and positivity are uniform on the
bounded range.  The maximum principle removes the homogeneous
Dirichlet kernel.  Linear Dirichlet solvability and Schauder estimates
therefore make $L:Y_R\to C^{1,\alpha}$ an isomorphism
\cite{GilbargTrudinger2001}.  The implicit function theorem proves openness in
$\beta$.  The bounds above prove closedness: subsequences converge in
lower H\"older norms, their limits solve the equation, and the uniform
$C^{3,\alpha}$ estimates retain that regularity.  At $\beta=0$ the
equation is $\Delta_gf=\eta_Nf$, with ordinary Dirichlet solvability.
Thus the scalar equation is solvable at $\beta=1$.

For uniqueness, at a positive maximum of the difference of two
solutions their gradients agree.  Their matrices and positive
prefactors therefore agree, whereas strict height monotonicity makes
the source difference positive, contradicting the Hessian comparison.
The reverse comparison completes uniqueness.  Applying the same
implicit function theorem with parameters $(s,Z)$ proves continuous
dependence into $C^{3,\alpha}$.  At stage junctions use one-sided
parameter neighborhoods and the agreement of the prescriptions.
The estimates were uniform on bounded trial sets, which proves the
last assertion.
\end{proof}

\subsection{The frozen return map and degree}

For each trial $Z$, the scalar problem has determined $f_s[Z]$;
the second equation remains to be imposed. We freeze its
coefficients and source at the trial pair and solve for a new
scalar function. A fixed point of this return map solves both
equations. Arbitrary trial inputs are not assumed to satisfy the floor.

For a trial pair $(f_s[Z],Z)$, evaluate all expressions in the
homotopy at that pair.  Write $A=u\Achi$, let
$B_s=\lambda u\Achi K(w,\cdot)$ be its drift, and write $Q_s$ and $q_s$
for the prescribed interior source and inner flux, including the
common terminal scale.  Define $T_sZ=Y$ by the linear problem
\begin{equation}\label{four:ex:frozen-problem}
 \begin{aligned}
 \divg_g(3A\grad Y+B_s)&=Q_s &&\hbox{in }\Omega_R,\\
 (3A\grad Y+B_s)\cdot\nu&=q_s &&\hbox{on }B,\\
 Y&=0 &&\hbox{on }S_R.
 \end{aligned}
\end{equation}
Thus only the occurrence of $3\grad Z$ in the flux is replaced by
$3\grad Y$; the coefficients, source, and boundary expression are
frozen.  In $q_s$ the value of $F$ is always its trace prescription.

The matrix $A$ is symmetric positive definite and $C^{1,\alpha}$.
On bounded trial sets its ellipticity constants, $C^{1,\alpha}$ norm,
and those of $B_s,q_s$ are bounded; $Q_s$ is bounded in $C^{0,\alpha}$.
In particular the source uses $\dd Z$ and $\Hess f$, whose regularity
is already available.  With
\[
 \mathcal H_R=\{Y\in H^1(\Omega_R):Y|_{S_R}=0\},
\]
the weak problem has coercive bilinear form
$3\int_{\Omega_R} A\grad Y\cdot\grad\varphi$.
Indeed $\Omega_R$ is connected and $S_R$ is nonempty, so Poincar\'e's
inequality holds on $\mathcal H_R$.  The prescribed flux is a bounded
linear functional by the trace inequality.  Lax--Milgram gives a unique
weak solution.  At an inner face $f$ is constant, so
$\Achi\nu=\chi\nu$, and its flux condition is equivalently
\[
 \partial_\nu Y=\frac{q_s-B_s\cdot\nu}{3u\chi}.
\]
This is a $C^{1,\alpha}$ ordinary normal datum.  The disjoint Dirichlet
and normal faces have no corners.  Linear boundary and interior
estimates give $Y\in C^{2,\alpha}$ with bounds uniform on bounded
trial sets; the weak energy estimate supplies the lower norm in these
estimates.  Coefficient dependence and uniqueness, or the same linear
estimates applied to differences, give continuous dependence.
Consequently
\begin{equation}\label{four:ex:compact-map}
 T:[0,4]\times X_R\longrightarrow X_R
 \quad\hbox{is continuous and compact on bounded sets}.
\end{equation}
Here compactness uses the compact embedding
$C^{2,\alpha}\hookrightarrow C^{1,\alpha}$ on the fixed truncation.

Every fixed point is smooth.  Initially it has $Z\in C^{2,\alpha}$
and $f\in C^{3,\alpha}$.  The trace equation raises the regularity of
$f$, and then expansion of the divergence equation, with its normal
boundary datum, raises the regularity of $Z$.  Repetition gives all
orders.  The smooth a priori estimates from the preceding sections
therefore apply to every fixed point.

\begin{proposition}[Existence and exhaustion]\label{four:prop:existence}
Fix the prepared geometry and $0<\eps<1$.  For every sufficiently large
$N$ and every allowed $R$, system~\eqref{four:eq:system} has a smooth
solution on $\overline\Omega_R$ with $t>-\eps$.  For each such fixed
$N$ there is a sequence $R\to\infty$ whose solutions converge smoothly
on compact subsets up to $B$ to a solution on $\overline\Omega$ with
$t\ge-\eps$.  This solution satisfies the prescribed inner boundary
conditions, the height interval~\eqref{four:eq:height-interval}, and the
weighted trace and flux estimates of Lemma~\ref{four:lem:weighted-trace}.
\end{proposition}

\begin{proof}
Choose $N$ large enough for Proposition~\ref{four:prop:height-exclusion},
the collar estimates, and the polynomial absorption conditions used
in Proposition~\ref{four:prop:bounded-range}.  Fix an allowed $R$.
By \eqref{four:eq:bounded-range} and
Proposition~\ref{four:prop:coupled-regularity}, choose $M_0$ strictly larger
than every $C^{1,\alpha}$ norm of an above-floor fixed point on this
truncation.  Define
\begin{equation}\label{four:ex:moving-sections}
 \mathcal O_s=\left\{Z\in X_R:
 \|Z\|_{C^{1,\alpha}}<M_0,\quad
 \min_{\overline\Omega_R}t\bigl(Z,\dd f_s[Z]\bigr)>-\eps\right\}.
\end{equation}
Lemma~\ref{four:ex:scalar-trial} and smooth implicit dependence of $t$ make
their union relatively open in $[0,4]\times X_R$.

Apply Lemma~\ref{b:lem:moving-degree} to $I=[0,4]$, $X=X_R$,
$T_s$ from \eqref{four:ex:frozen-problem}, and
$\mathcal U=\{(s,Z):Z\in\mathcal O_s\}$.
The norm cutoff makes $\mathcal U$ bounded; the scalar solution map
makes it relatively open; and \eqref{four:ex:compact-map} supplies
continuity and compactness. Every fixed point is smooth. A fixed
point on its relative boundary would satisfy $t\ge-\eps$ and meet
either the norm cutoff or the floor. The first is excluded by the
choice of $M_0$, and the second by
Proposition~\ref{four:prop:floor-exclusion}. The section degree is
therefore constant. This checks the varying-domain hypotheses for
this boundary problem without imposing the floor on trial inputs.

At $s=4$, the scalar prescription is
$\tr_{\Achi}H^f=h+D_1(x,w)$ with zero boundary heights.
At an interior extremum $w=0$ and $D_1=0$, so the maximum principle
gives $f=0$ for every trial $Z$, and then $t=Z$.  The terminal scale
sets both frozen right sides to zero and the drift is absent.
Testing the frozen equation with $Y$ gives $Y=0$; thus $T_4$ is the
zero map, even though its positive coefficients may depend on the
input.  Its sole fixed point, $Z=0$, belongs to $\mathcal O_4$ and
has degree one.  The degree at $s=0$ is consequently one, proving
existence there.

Now keep $N$ fixed.  The preceding choices of sufficiently large $N$
used only the already proved uniform height and collar results and
polynomial losses such as $\Pi_N\eta_N/\ell\to0$.  The arbitrary
fixed-$N$ constants in the scalar construction impose no subsequent
condition on $N$.  The bounds in Proposition~\ref{four:prop:coupled-regularity}
are uniform in allowed $R$ on each compact set, including its portions
on $B$.  A diagonal Arzel\`a--Ascoli argument along $R\to\infty$ gives
smooth convergence on all these sets.  The equation and inner data
pass to the limit, as do $t\ge-\eps$, the height interval, and the
collar estimates.  Smooth convergence on the fixed collars passes
the weighted trace and flux inequalities as well.  Strictness of the
floor in this exhaustion limit is unnecessary below.
\end{proof}

\subsection{Decay on the end}

\begin{lemma}[End estimates]\label{four:lem:end-decay}
For every solution obtained in Proposition~\ref{four:prop:existence}, with
$N$ fixed, there is $a_N>0$ such that, for every integer $j\ge0$,
\begin{equation}\label{four:ex:end-f}
 |\nabla^jf|+|\nabla^j(t-Z)|\le C_j(N)e^{-a_Nr}
 \quad\hbox{on the distant end}.
\end{equation}
Moreover
\begin{equation}\label{four:ex:end-tz}
 t,Z=O_2(r^{-2}),\qquad \Delta_gt=O(r^{-4-\delta_1}).
\end{equation}
The constants in these assertions may depend arbitrarily on $N$.
The zeroth-order estimates and \eqref{four:ex:end-f} hold uniformly for the
finite-truncation solutions up to their outer faces.
\end{lemma}

\begin{proof}
Choose a fixed sphere beyond the supports of $K,C$, and all collar
terms.  Along each finite-truncation solution the trace equation is
the homogeneous linear equation
\begin{equation}\label{four:ex:f-end-linear}
 \Achi:\Hess_g f-c(x)f=0,\qquad
 c(x)=\eta_N\frac{\sqrt D}{l}\ge c_N>0.
\end{equation}
By \eqref{four:eq:bounded-range} the coefficients are uniformly elliptic;
by Proposition~\ref{four:prop:coupled-regularity} all their derivatives are
bounded on uniform unit patches, independently of $R$.  On the end,
\[
 (\Achi:\Hess_g-c)e^{-ar}
 =e^{-ar}\bigl(a^2\Achi(\dd r,\dd r)
          -a\Achi:\Hess_g r-c\bigr).
\]
Choose $a=a_N>0$ so small that the first term is less than $c_N/4$,
and then enlarge the fixed sphere so the second term has magnitude
less than $c_N/4$.  Thus $e^{-a_Nr}$ is a strict supersolution.
A fixed multiple dominates both signs of $f$ on the initial sphere
and dominates the zero Dirichlet data on $S_R$.  Linear comparison
gives $|f|\le C(N)e^{-a_Nr}$ uniformly in $R$.  Local estimates for
\eqref{four:ex:f-end-linear}, including the zero-Dirichlet estimates on the
uniform outer-sphere charts, give the same exponential decay for
every derivative, after harmless reduction of $a_N$ if necessary.
The identity
\[
 Z-t=\tfrac16\log(1+l(t)^2|\dd f|^2)
\]
and the uniform smooth local bounds then give
\eqref{four:ex:end-f}.  Every graph contribution to the metric and to the
equations is exponentially small with unit-patch derivatives.

At zero graph and $K=0$, the exact formulas give
\[
 u=e^{2Z}l(Z),\qquad V=3u\grad Z,
 \qquad \cT=3(p(Z)+1)|\dd Z|^2.
\]
Hence the actual divergence equation, divided by $u$, becomes
\begin{equation}\label{four:ex:Z-end-equation}
 3\Delta_gZ+3\mathfrak b_N(Z)|\dd Z|_g^2
 =\rho-m_0(Z)C_N\rho_0+\mathcal E_N,
 \quad \mathfrak b_N(z)=p(z)+2-\delta_0(p(z)+1),
\end{equation}
where $\mathcal E_N$ and its derivatives on unit patches are
$O_j(e^{-a_Nr})$.  Define the strictly increasing function
\begin{equation}\label{four:ex:Phi}
 \Phi(0)=0,\qquad
 \Phi'(z)=\exp\left(\int_0^z \mathfrak b_N(s)\,\dd s\right).
\end{equation}
Its derivative is bounded above and away from zero on the fixed
bounded range of $Z$.  The chain rule cancels the quadratic term:
\begin{equation}\label{four:ex:Phi-Poisson}
 \Delta_g\Phi(Z)
 =\frac{\Phi'(Z)}3
    \bigl(\rho-m_0(Z)C_N\rho_0+\mathcal E_N\bigr),
 \qquad |\Delta_g\Phi(Z)|\le C(N)r^{-4-\delta_1}.
\end{equation}
Put $\beta_1=\delta_1/2$ and
$W(r)=r^{-2}(1-r^{-\beta_1})$.  In dimension four,
\[
 \Delta_\delta W=-\beta_1(2+\beta_1)r^{-4-\beta_1},
 \qquad
 \Delta_gW=-\beta_1(2+\beta_1)r^{-4-\beta_1}+O(r^{-6}).
\]
Since $0<\beta_1<\delta_1<1$, a sufficiently large multiple of $W$
is positive and has negative Laplacian dominating the absolute source
in \eqref{four:ex:Phi-Poisson} on a sufficiently distant end.  Enlarge the
multiple to dominate $|\Phi(Z)|$ on its initial sphere.  Its value is
positive on $S_R$, whereas $\Phi(Z)=0$ there.  Applying the maximum
principle to both $\Phi(Z)-CW$ and $-\Phi(Z)-CW$ gives
\[
 |\Phi(Z)|\le C(N)r^{-2},\qquad |Z|\le C(N)r^{-2},
\]
uniformly on truncations and then in the exhaustion limit.

Here are the differentiated estimates, which do not follow from the
pointwise barrier alone.  On an annulus of radius $r_1$, rescale by
$x=r_1y$ and set $U(y)=r_1^2\Phi(Z(r_1y))$.
Its supremum and its Poisson source in the rescaled equation are
bounded: the source is $r_1^4$ times the right side of
\eqref{four:ex:Phi-Poisson}, of size $O_N(r_1^{-\delta_1})$.
The rescaled background coefficients have uniform smooth bounds.
Interior $W^{2,p}$ estimates, for a fixed $p>4$, give bounded
$C^{1,\theta}$ norms of $U$ on a smaller annulus.  The weights have
their differentiated power decay, the compositions with $Z$ now
have bounded scaled H\"older norms, and the exponential errors remain
bounded with all scaled derivatives.  The right side therefore has a
bounded $C^{0,\theta}$ norm.  Schauder estimates give a bounded
$C^{2,\theta}$ norm of $U$.  The inverse of $\Phi$ has bounded smooth
derivatives on the relevant range.  Rescaling back gives
$Z=O_2(r^{-2})$, and \eqref{four:ex:end-f} gives the same statement for $t$.
Finally $|\dd Z|^2=O(r^{-6})$ in \eqref{four:ex:Z-end-equation}; since
$\delta_1<1$, this is lower order than $r^{-4-\delta_1}$.
Equation~\eqref{four:ex:end-f} now gives the asserted bound for $\Delta_gt$.
\end{proof}

\subsection{Scalar curvature, boundary sign, and finite flux}

For the desired system $F=h+C$ and
$w(h)=\eta_N a\sigma$.  Substitute these into
\eqref{four:eq:scalar-identity} and use $\divg V=u\Xi$ to obtain
\begin{equation}\label{four:ex:positive-scalar}
 \begin{split}
 \tfrac12e^{2t}\Scal_{\ghat}
 ={}&\mu+J(w)+w(C)-\rho-(h+C)\tau\\
 &+(1-\delta_0)(\cT+\eta_N a\sigma)+m_0\cP\ \ge0.
 \end{split}
\end{equation}
Indeed $|w|\le1$, and on $\supp K$ the height interval
\eqref{four:eq:height-interval} permits use of \eqref{four:eq:trace-slack}:
\[
 \mu+J(w)+w(C)-\rho-(h+C)\tau
 \ge\mu-|J|-|\dd C|-\rho-|(h+C)\tau|\ge0.
\]
Outside $\supp K$ the same slack inequality needs no height restriction.
The other terms are nonnegative by \eqref{four:eq:coercivity} and the
definitions of the penalty and $\delta_0$.  Thus the nonmaximal trace
term has been retained and controlled at this final step.

The inner flux condition and \eqref{four:eq:boundary-identity} give exactly
\begin{equation}\label{four:ex:negative-boundary}
 H+3\partial_\nu t=-N_B,
 \qquad \widehat H=-e^{-t}\sqrt d\,N_B<0
 \quad\hbox{on every component of }B.
\end{equation}
All quantities here are finite and $d>0$ for each fixed $N$.
Furthermore \eqref{four:ex:end-f} and \eqref{four:ex:end-tz} imply
\begin{equation}\label{four:ex:final-AF}
 \ghat-\delta=O_2(r^{-2}),\qquad
 \Scal_{\ghat}=O(r^{-4-\delta_1}).
\end{equation}
For the scalar estimate one may use the conformal scalar formula
with $\gbar=g+O_j(e^{-a_Nr})$:
$e^{2t}\Scal_{\ghat}=\Scal_{\gbar}-6\Delta_{\gbar}t
-6|\dd t|_{\gbar}^2$.  The prepared scalar decay and
\eqref{four:ex:end-tz} give precisely \eqref{four:ex:final-AF}.
The nonnegative scalar curvature is therefore integrable.
The inequality
\begin{equation}\label{four:ex:metric-comparison}
 \ghat=e^{2t}(g+l^2\dd f\otimes\dd f)\ge e^{-2\eps}g
\end{equation}
also proves completeness with $B$ included: a $\ghat$-Cauchy sequence
is $g$-Cauchy, its limit belongs to the closed complete exterior, and
smooth local metric equivalence gives convergence in $\ghat$.

\begin{proposition}[Energy comparison]\label{four:prop:mass-comparison}
For each fixed $\eps>0$ and all sufficiently large $N$, the limiting
metric has finite ADM energy and
\begin{equation}\label{four:ex:energy-error}
 E_{\ghat}\le E_g+\frac{\Pi_N}{3\omegaThree}
                  \left(\ell+\frac{\ell^2}{\eta_N}\right).
\end{equation}
The polynomial $\Pi_N$ depends only on the prepared data, $\eps$,
and the fixed geometric choices.  In particular the error tends to
zero as $N\to\infty$ with $\eps$ fixed.
\end{proposition}

\begin{proof}
We first establish the exact sign and normalization at fixed $N$.
Lemma~\ref{four:lem:end-decay} gives $u=1+O_N(r^{-2})$ and
\begin{equation}\label{four:ex:V-asymptotics}
 V=3\grad_gt+O_N(r^{-5})+O_N(e^{-a_Nr}).
\end{equation}
Also $u\Xi\in L^1(\Omega,\dd V_g)$: on the end the quadratic term
is $O_N(r^{-6})$, the graph and height-gradient terms are exponentially
small, and the remaining weights are $O_N(r^{-4-\delta_1})$.
The divergence theorem on a large truncation consequently gives a
finite limit
\begin{equation}\label{four:ex:flux}
 \mathfrak F:=\lim_{r\to\infty}\int_{S_r}V\cdot\nu_g\,\dd A_g
   =\int_\Omega u\Xi\,\dd V_g+\int_B V_\nu\,\dd A_g.
\end{equation}
The plus sign before the inner flux arises because the domain-outward
normal on $B$ is $-\nu$.

To compute the metric flux, write $\ghat=e^{2t}g$ up to exponential
errors.  The leading perturbation relative to $g$ is $2t\delta_{ij}$,
so in four dimensions the change of its energy numerator is
\[
 \partial_j(2t\delta_{ij})-\partial_i(2t\delta_{jj})
 =-6\partial_i t.
\]
All omitted derivatives are $O_N(r^{-5})$ or exponentially small.
Their sphere integrals tend to zero.  Since $\dd t=O_N(r^{-3})$ and
$g-\delta=O_2(r^{-2})$, changing the normal and measure in its flux
from Euclidean to background ones also has vanishing error.  Equation
\eqref{four:ex:V-asymptotics} thus proves both existence of the new energy
and the exact relation
\begin{equation}\label{four:ex:mass-flux}
 E_{\ghat}-E_g
 =-\frac1{\omegaThree}\lim_{r\to\infty}
       \int_{S_r}\partial_{\nu_\delta}t\,\dd A_\delta
 =-\frac{\mathfrak F}{3\omegaThree}.
\end{equation}

We now use polynomial constants only.  Apply
Lemma~\ref{four:lem:weighted-trace} with $\varphi=1$ and use
\eqref{four:eq:boundary-flux-bound}.  On the fixed compact collars
$\mathcal C$ one has $\min_{\mathcal C}\rho>0$, and therefore
\[
 \begin{split}
 \int_B|V_\nu|\,\dd A_g
 &\le\Pi_N\frac{\eta_N}{\ell}
       \int_{\mathcal C}u(1+\sqrt{\cT})\,\dd V_g\\
 &\le\Pi_N\frac{\eta_N}{\ell}
       \int_{\mathcal C}u(\rho+\delta_0\cT)\,\dd V_g.
 \end{split}
\]
The polynomial absorbs the fixed comparison constant in the second
line.  Because $\eta_N/\ell=e^{-\eps N/2}$, for sufficiently large
$N$ this is at most one half of
$\int_\Omega u(\rho+\delta_0\cT)$.

It remains to estimate the penalty.  Its support in an above-floor
solution lies in
$-\eps\le t\le-\eps+1/N$.  For $N>1/\eps$ this interval is negative,
and uniformly there
\begin{equation}\label{four:ex:penalty-band}
 l\asymp\ell,\qquad L_0\asymp\ell,\qquad
 u\le C(\ell+\ell^2\sigma),\qquad
 uv\le C\ell^2\sigma,\qquad
 ua\sigma=L_0l\sigma^2\ge c\ell^2\sigma^2.
\end{equation}
For example $uv=L_0l^2\sigma^2/\sqrt{1+l^2\sigma^2}
\le L_0l\sigma$.  Thus on that band
\[
 um_0\cP\le\Pi_N\bigl[\ell\rho_0+
                         \ell^2\sigma(\rho_0+\rho_1)\bigr].
\]
Young's inequality applied to
$\sqrt{\eta_N}\ell\sigma$ and
$\Pi_N\ell(\rho_0+\rho_1)/\sqrt{\eta_N}$, together with
\eqref{four:ex:penalty-band}, yields everywhere
\begin{equation}\label{four:ex:penalty-absorption}
 um_0\cP\le\tfrac12\delta_0u\eta_N a\sigma
   +\Pi_N\left[\ell\rho_0+
       \frac{\ell^2}{\eta_N}(\rho_0^2+\rho_1^2)\right].
\end{equation}
The enlarged $\Pi_N$ is still polynomial.  The three residual weights
are integrable against $\dd V_g$ in dimension four:
$\rho_0=r^{-4-\delta_1}$, $\rho_0^2=r^{-8-2\delta_1}$, and
$\rho_1^2=r^{-6}$ on the end, and they are bounded on the compact part.
Combining \eqref{four:ex:flux}, the boundary absorption, and
\eqref{four:ex:penalty-absorption} gives
\[
 \mathfrak F\ge-\Pi_N\left(\ell+\frac{\ell^2}{\eta_N}\right).
\]
This and \eqref{four:ex:mass-flux} prove \eqref{four:ex:energy-error}.
Finally $\ell=e^{-\eps N}$ and $\ell^2/\eta_N=e^{-\eps N/2}$;
both dominate every polynomial loss.  No constant from the fixed-$N$
decay or Schauder estimates appears in this bound.
\end{proof}


\subsection{Completion of the boundary deformation}

\begin{proof}[Proof of Theorem~\ref{four:thm:boundary-deformation}]
The geometric preparation in
Proposition~\ref{four:prop:geometric-preparation} chooses a closed connected
one-ended exterior $\Omega$ inside the original data. It has nonempty
compact smooth boundary and does not change the asymptotic energy.
The last assertion of that Proposition compares every full cut in $\Omega$
with the original infimum $A_*$. Fix $0<\eps<1$. For each sufficiently
large $N$, Proposition~\ref{four:prop:existence} gives the desired smooth
solution after exhausting the outer radius at that fixed $N$.
Equation~\eqref{four:ex:positive-scalar} gives nonnegative scalar curvature;
Equation~\eqref{four:ex:negative-boundary} gives strictly negative inner
mean curvature on every face. Equations~\eqref{four:ex:final-AF} and
\eqref{four:ex:metric-comparison} establish the required asymptotics,
integrability and completeness, as well as the lower metric bound.
Proposition~\ref{four:prop:mass-comparison} gives finite ADM energy with
the stated error. Its polynomial bound uses only the fixed geometric
data and $\eps$. Since
$\ell=e^{-\eps N}$ and $\ell^2/\eta_N=e^{-\eps N/2}$, this error
tends to zero. No estimates uniform over a further sequence of prepared
data, and no convergence of the constructed metrics as $N\to\infty$,
are needed.
\end{proof}
